%----------
% DISCUSIÓN — lectura transversal de los resultados
%----------

Más allá de cada tabla, los resultados del capítulo admiten una lectura de
conjunto en torno a tres ideas: el carácter condicional de las \textit{Agent
Skills}, la lección metodológica que deja la evaluación pormenorizada que hemos hecho, y el techo que impone la
propia tarea, del que se derivan las implicaciones para el proyecto.

\subsection{El valor de las \textit{Agent Skills} es condicional al modelo}

El hallazgo central obliga a matizar una expectativa extendida: que dotar a un modelo
de más componentes (herramientas, conocimiento cargable bajo demanda, capacidad de
actuar) mejora por sí solo su desempeño. No es lo que se observa. Empaquetar la
metodología experta como \textit{skills} cargables no eleva el acuerdo de forma
automática, y lo que separa a los cuatro modelos no es tanto la dirección del cambio
como su magnitud. Solo \texttt{gemini} registra una ganancia de entidad, la de
\texttt{masc\_generico}, donde su $\kappa$ sube 0,188 puntos. En los demás los
movimientos se quedan en unas pocas milésimas o van hacia abajo, y
\texttt{gpt-5.4-nano} no mejora en ninguna de las cinco variables. La arquitectura no
es útil ni inútil en abstracto, sino que \textbf{amplifica lo que cada modelo ya
trae}. Aprovechar la carga bajo demanda exige que el modelo sepa decidir qué necesita
y cuándo, una competencia desigual entre modelos, y cuando falta, el componente aporta
contexto que satura en lugar de ayudar.

De la descomposición por componentes (Sección~\ref{subsec:iris_descomposicion}) se
sigue además una recomendación práctica, con la cautela de que ninguna configuración
se impone en las cinco variables. La recuperación en vivo sobre las guías no produce
una ganancia que justifique su coste, ya que sobre la configuración mínima mejora el
F1 de la clase positiva en nueve de las veinte celdas y lo empeora en ocho. La
preferible para un despliegue real es la de \textit{skills} con sus resúmenes, sin
herramienta, que encabeza siete de las veinte celdas, más que ninguna otra, y es a la
vez la más simple y la más barata. No es, con todo, una opción imbatible, ya que queda
por detrás de la configuración mínima en siete celdas, y en una de ellas por mucho, la
de \texttt{gpt-4o-mini} en \texttt{masc\_generico}, donde su F1 cae de 0,765 a
0,401.


\subsection{Exactitud, \texorpdfstring{$\kappa$}{kappa} y prevalencia: por qué informar de las tres}
\label{subsec:iris_tres_metricas}

El segundo aprendizaje es metodológico. En tareas con clases muy desequilibradas, la
exactitud y el coeficiente $\kappa$ pueden desacoplarse y llegar a moverse en
sentidos opuestos. Una intervención que vuelve al modelo menos conservador aumenta
los aciertos triviales por efecto de la distribución marginal, sin que mejore su capacidad
discriminativa. Este trabajo lo ilustra en dos celdas opuestas. En
\texttt{lenguaje\_sexista}, cuya clase mayoritaria es el \textit{Sí},
\texttt{gpt-4o-mini} se vuelve menos conservador al recibir las \textit{skills} y pasa
de marcar el 9,0\,\% de las piezas al 46,2\,\%, con lo que su exactitud sube de 0,236 a
0,457 mientras su $\kappa$ baja de $-0,013$ a $-0,037$ (Tabla~\ref{tab:var-v25}). En
\texttt{asimetria\_mujer\_hombre}, cuya clase mayoritaria es el \textit{No},
\texttt{gpt-5.4-nano} hace justo lo contrario y pasa del 37,5\,\% al 0,7\,\%, de modo
que su exactitud salta de 0,619 a 0,931 al tiempo que su recall cae de 0,457 a cero
(Tabla~\ref{tab:var-v33}). En ninguno de los dos casos discrimina mejor el modelo, sino
que se acerca a la respuesta que más se repite y la exactitud premia ese acercamiento.
Apoyarse solo en ella llevaría a concluir que las \textit{Agent Skills} mejoran el
sistema, cuando el $\kappa$ indica lo contrario. Por eso, en tareas
desequilibradas, la exactitud debe interpretarse siempre junto al $\kappa$ y a la
prevalencia, para no tomar un cambio en la distribución de respuestas por una mejora
real de criterio.


\subsection{Dos límites, uno de la tarea y otro de los modelos}
\label{subsec:iris_dos_limites}

El acuerdo absoluto es bajo, y el análisis del Capítulo permite repartir ese
desajuste entre dos causas que conviene no confundir.

La primera es un \textbf{límite de la tarea}. Dos equipos expertos codificaron la
misma variable con criterios muy distantes, hasta el punto de que
\texttt{sexismo\_discurso} pasa del 11,5\,\% al 65,2\,\% de piezas marcadas según
quién anote (Sección~\ref{subsec:iris_errores}). La brecha, además, no es
institucional. Dentro de un mismo equipo esa variable va del 28,3\,\% al 97,8\,\%
según quién codifique, un rango de 69,5 puntos que por sí solo supera a los 53,7 que
separan a los dos equipos, de modo que dos personas con la misma formación y el mismo
libro de códigos marcan, una, poco más de una de cada cuatro piezas, y la otra, casi
todas. Cuando el criterio admite esa
amplitud, el patrón de referencia deja de ser un suelo estable y el coeficiente
$\kappa$ mide en buena parte distancia entre criterios más que error del sistema.
Esta variabilidad no es un defecto del corpus de IRIS, sino un rasgo reconocido de
las tareas de anotación subjetiva, que la investigación reciente propone tratar
como información y no como ruido a eliminar~\cite{davani2022, plank2022}. La
campaña EXIST, referencia en identificación de sexismo en español, ha adoptado por
ello un paradigma de aprendizaje con desacuerdo en el que se conservan las
etiquetas de todas las personas anotadoras en lugar de agregarlas en una
sola~\cite{plaza2024}.

La segunda es un \textbf{límite de los modelos}, y sería un error disolverlo en el
anterior. En \texttt{lenguaje\_sexista} ambos equipos coinciden en marcar la
variable en la mayoría de las piezas, el 74,6\,\% y el 85,5\,\% respectivamente, y
aun así tres de los cuatro modelos presentan un recall prácticamente nulo. Ahí no
hay criterio en disputa, hay un fenómeno que las personas ven y el sistema no. El
caso examinado en la Sección~\ref{subsec:iris_caso} muestra el mecanismo, ya que
los modelos localizan el masculino genérico y lo declaran un colectivo mixto
válido, reproduciendo el mismo automatismo que se les pedía detectar.

El peso de cada límite cambia según la variable, y de ahí se sigue que el margen de
mejora tenga también dos vías. Una consiste en explicitar mejor el constructo en el
libro de códigos, apoyándose en anotaciones cualitativas que lo delimiten con
ejemplos, y conviene priorizarla en las variables de criterio más implícito. La
otra consiste en mejorar la sensibilidad del sistema en aquello que las expertas sí
identifican de forma concordante. Conviene además evitar una simplificación, ya que
los modelos no comparten un sesgo homogéneo, sino cuatro comportamientos distintos
que solo coinciden en su dirección.

\subsection{Implicaciones para IRIS}
\label{subsec:iris_implicaciones}

De los resultados se desprenden cuatro orientaciones prácticas para el proyecto.

\paragraph{El despliegue local es viable.} El modelo abierto \texttt{gemma},
ejecutado en infraestructura propia, compite con los servicios comerciales y los
supera en algunas variables, como \texttt{asimetria\_mujer\_hombre}, donde obtiene
el mejor recall y la mejor precisión de los cuatro
(Sección~\ref{subsec:iris_local}). Para una institución que maneja material sujeto
a privacidad, esto significa que preservar los datos y prescindir de tarifas no
obliga a renunciar al rendimiento, a cambio de asumir un coste de cómputo. Ese coste,
además, se comporta al revés que el de los modelos de pago, en los que el empaquetado
multiplica la factura por 2,2. En el local la reduce, ya que procesar una pieza pasa
de 319,3 segundos en el nivel de control a 218,6 en el desplegado, porque la carga bajo
demanda incorpora solo la \textit{skill} que hace falta en lugar de arrastrar la
metodología completa en cada llamada.

\paragraph{Conviene tratar cada variable por separado.} Ni la configuración óptima
ni la utilidad del sistema son las mismas en las cinco. En
\texttt{masc\_generico} el sistema detecta ya casi tres de cada cuatro casos
señalados por las expertas y puede incorporarse al flujo de trabajo, mientras que en
\texttt{lenguaje\_sexista} no aporta información utilizable pese a ser la variable
que más presupuesto de inferencia consume, el 24,2\,\% del total
(Sección~\ref{subsec:iris_coste}).
Concentrar el esfuerzo donde el sistema funciona, y no distribuirlo por igual,
mejora el rendimiento y reduce el gasto a la vez. Esa afirmación admite una cifra.
Asignar a cada variable el modelo que mejor se comporta en ella supone 8,77 USD sobre
las 1.313 piezas, cantidad en la que \texttt{asimetria\_mujer\_hombre} no devenga
tarifa por recaer en el modelo local, si bien consume tiempo de cómputo. Fijar en
cambio un único modelo para las cinco costaría 13,16 USD con \texttt{gemini}, que es el
de mayor acuerdo, y 6,89 USD con \texttt{gpt-4o-mini}, que sale más barato pero queda
por detrás en tres de las cinco variables. La asignación por variable se sitúa, por
tanto, entre ambas en coste y por delante de las dos en rendimiento, que es cuanto cabe
pedirle a una decisión de despliegue.

\paragraph{La supervisión humana no es un remedio provisional.} Que el sistema
requiera revisión experta no constituye una anomalía de esta implementación, sino
lo que la literatura viene recomendando para las tareas de anotación con carga
interpretativa. Pangakis et al.~\cite{pangakis2023} replican veintisiete tareas de
anotación procedentes de la investigación en ciencias sociales y concluyen que
ningún proceso automático puede darse por válido sin contrastarlo con etiquetas
humanas, ya que en nueve de ellas la precisión o el recall caen por debajo de 0,5.
Kocoń et al.~\cite{kocon2023} constatan que el rendimiento de un modelo general se
degrada de forma acusada justo en las tareas subjetivas, con una pérdida media del
25,5\,\% frente a los sistemas especializados. Y desde la perspectiva de la
gobernanza, Gorwa et al.~\cite{gorwa2020} advierten de que los sistemas
algorítmicos de moderación, aun bien optimizados, siguen siendo opacos y difíciles
de fiscalizar. El paradigma \textit{human-in-the-loop} que adopta este trabajo se
sitúa, por tanto, en el consenso de la disciplina.

\paragraph{Ante la duda, conviene marcar de más.} Como herramienta de apoyo a la
postedición, el sistema no decide, sino que selecciona lo que merece una segunda
lectura. En ese uso los dos errores no cuestan lo mismo, ya que un falso positivo
supone unos segundos de descarte para la analista mientras que un falso negativo
es un caso que nadie revisará (Sección~\ref{sec:metricas_eval}). Los resultados
apuntan en la dirección contraria a la deseable, puesto que el sistema falla
precisamente por defecto, de modo que cualquier ajuste futuro debería orientarse a
elevar el recall aun a costa de la precisión. La Sección~\ref{subsec:iris_equipo_ia}
apunta una vía concreta en ese sentido, la de marcar la pieza cuando lo haga
cualquiera de los modelos, que eleva el recall en las cinco variables, si bien obliga
a ejecutarlos todos sobre cada pieza.

Con estas cuatro orientaciones, el sistema encaja mejor como asistente que como
decisor. Su salida trazable, con código, explicación y evidencias literales,
permite que la analista audite, acepte o corrija cada predicción, y el caso
analizado en la Sección~\ref{subsec:iris_caso} muestra que esa explicación no es un
añadido cosmético, sino lo que permite entender por qué falla el sistema y, con
ello, corregirlo. El valor no está en sustituir el criterio experto, sino en
ordenar el corpus, señalar los casos que merecen atención y hacerlo de forma
auditable.
